\section{Biological Example Family I: Sleep-Wake Switching}

The strongest biological transition model available to this manuscript remains
the sleep-transition source from Chapter~\ref{ch:consciousness-biology}. It
matters here because it supports a disciplined use of bifurcation and threshold
language.

\begin{discussion}
Sleep-wake switching is pedagogically valuable for three reasons:
\begin{enumerate}[nosep]
  \item it is a genuine biological regime transition rather than a self-help
        anecdote;
  \item it demonstrates that state changes can be abrupt, patterned, and
        parameter-sensitive;
  \item it shows that the key mathematical question is often not ``how much
        effort is present?'' but ``which stability condition is being changed?''
\end{enumerate}
\end{discussion}

This does \emph{not} imply that reading a theorem and falling asleep are the
same mechanism. The evidence lane remains explicit. The strong lane supports
the general legitimacy of regime-switching language in biology. The transfer to
behavioural contexts remains a \textbf{medium} interpretive move. That move is
still useful, because it encourages the reader to think in terms of thresholds,
timing, and basin geometry rather than mere desire.
